% Revize 2026-10-28: obsah, matematika a návaznost; podrobnosti a časy
% viz planovani/REVIZE_PRIPRAV_1_ROCNIK.md.
% Aktuální porovnání se skutečným obsahem místních PDF Realisticky.cz:
% externi_zdroje/realisticky/01_Zakladni_poznatky/09_Lomene_vyrazy/010905_Zjednodusovani_lomenych_vyrazu_I/010905_Lekce.pdf
% Revize 2026-10-27: cílené doplnění; nové úlohy a řešení jsou vlastní tvorba.
% Metodický vzor: 4_rocnik/pripravy/Uvod_do_kombinatoriky.pdf (samostatný dokument).
% Hodina 55, 1. ročník SPŠ. Rozsah: Rozpis_vyuky_MAT_1_rocnik.tex, řádek 55.

% Plán: 1_rocnik/tematicke_plany/SPS.pdf; společná projektová šablona.

% Učebnice: ucebnice/01_Zakladni_poznatky/Algebraicke_vyrazy/Algebraické výrazy.tex

% Původ úloh: P1: učebnice, „Umocňování a úpravy složených lomených výrazů“, přesný složený zlomek. Ostatní vlastní.

% Metodický zdroj: externi_zdroje/realisticky/01_Zakladni_poznatky/09_Lomene_vyrazy/010905_Zjednodusovani_lomenych_vyrazu_I/010905_Lekce.pdf

% URL: http://www.realisticky.cz/ucebnice/01%20Matematika%20S%C5%A0/01%20Z%C3%A1kladn%C3%AD%20poznatky/09%20Lomen%C3%A9%20v%C3%BDrazy/05%20Zjednodu%C5%A1ov%C3%A1n%C3%AD%20lomen%C3%BDch%20v%C3%BDraz%C5%AF%20I.pdf

\documentclass[11pt,a4paper]{article}

% Spolecny styl pro ucitelske pripravy MAT (Work/MAT)
% Verze 1.1, 2026-10-23
\NeedsTeXFormat{LaTeX2e}
\ProvidesPackage{priprava}[2026/10/23 v1.1 Sablona ucitelskych priprav MAT]

% Zakladni typografie a cestina
\RequirePackage[T1]{fontenc}
\RequirePackage[utf8]{inputenc}
\RequirePackage[czech]{babel}
\RequirePackage{lmodern}
\RequirePackage{microtype}
\RequirePackage{amsmath,amssymb}
\RequirePackage{xcolor}
\RequirePackage{enumitem}
\RequirePackage[most]{tcolorbox}
\RequirePackage{tabularx,booktabs,array}
\RequirePackage{fancyhdr}
\RequirePackage{needspace}
\RequirePackage{geometry}
\RequirePackage[hidelinks]{hyperref}

\geometry{a4paper,top=17mm,bottom=18mm,left=23mm,right=23mm}
\definecolor{notegray}{gray}{0.39}
\setlength{\parindent}{0pt}
\setlength{\parskip}{2.5pt plus 1pt minus 0.5pt}
\setlist[itemize]{topsep=2pt,itemsep=1pt,parsep=0pt,leftmargin=1.45em}
\setlist[enumerate]{topsep=2pt,itemsep=2pt,parsep=0pt,leftmargin=1.65em}
\setcounter{secnumdepth}{0}
\renewcommand{\arraystretch}{1.12}

% Zahlavi: vlevo rocnik | globalni tema, vpravo cislo: nazev hodiny
% Zapati zustava zcela prazdne.
\pagestyle{fancy}
\fancyhf{}
\renewcommand{\headrulewidth}{0.35pt}
\renewcommand{\footrulewidth}{0pt}
\setlength{\headheight}{14pt}
\newcommand{\nastavPripravu}[4]{%
  \fancyhead[L]{\small #1\enspace\textbar\enspace #2}%
  \fancyhead[R]{\small #3: #4}%
}

% Minimalisticke cernobile boxy. Kazda definice zvlast.
\tcbset{commonlesson/.style={enhanced,breakable=false,colback=white,colframe=black,
  colbacktitle=white,coltitle=black,boxrule=0.45pt,
  arc=0.8mm,outer arc=0.8mm,left=3mm,right=3mm,top=1.6mm,bottom=1.6mm,
  toptitle=1mm,bottomtitle=0.7mm,fonttitle=\bfseries,
  attach boxed title to top left={xshift=2.5mm,yshift=-1.15mm},
  boxed title style={colback=white,colframe=white,boxrule=0pt,left=1mm,right=1mm}}}
\newtcolorbox{zadaniBox}[1][]{commonlesson,title={Zadání},#1}
\newtcolorbox{prikladBox}[1][]{commonlesson,title={Příklad},#1}
\newtcolorbox{definiceBox}[1][]{commonlesson,title={Definice},#1}

% Spolecny box pro zadani a jeho reseni, bez mezery mezi castmi.
% Pouziti:
%   \begin{prikladReseniBox}
%     Zadani ulohy.
%     \tcblower
%     \begin{align*} ... \end{align*}
%   \end{prikladReseniBox}
% Spodni cast ma jemny sedy podklad a delici linku; lze v ni sazet
% vice radku matematiky. Obe casti maji jen jeden spolecny ramecek.
\tcbset{solutionlesson/.style={
  commonlesson,
  enhanced jigsaw,
  breakable,
  colbacklower=black!3,
  segmentation style={solid,draw=black!25,line width=0.35pt},
  fontlower=\small\color{notegray},
  before lower={\noindent\textbf{Řešení:}\enspace},
  leftlower=3mm,rightlower=3mm
}}
\newtcolorbox{prikladReseniBox}[1][]{solutionlesson,title={Příklad},#1}
\newtcolorbox{zadaniReseniBox}[1][]{solutionlesson,title={Zadání},#1}

% Pomucky pro autora pripravy: sedive poznamky; \reseni pro starsi pripravy
\newcommand{\poznamka}[1]{\par\begingroup\small\color{notegray}\itshape
  \textnormal{\textbf{Poznámka pro učitele.}}\enspace #1\par\endgroup}
\newcommand{\reseni}[1]{\par\begingroup\small\color{notegray}
  \textbf{Očekávané řešení:}\enspace #1\par\endgroup}
\newcommand{\blok}[1]{\needspace{5\baselineskip}\subsection*{#1}}
\newcommand{\vyrok}[1]{\uv{#1}}


\nastavPripravu{1. ročník}{Algebraické výrazy}{55}{Složené lomené výrazy}

\begin{document}

\blok{Výklad a souvislosti}

Hlavní zlomková čára znamená dělení celého horního výrazu celým dolním. Podmínky proto pocházejí jak z vnitřních jmenovatelů, tak z požadavku, aby celý dolní výraz nebyl nulový.

\begin{definiceBox}[title={Definice},breakable=false]
\textbf{Složený zlomek} je zlomek, v jehož čitateli nebo jmenovateli vystupují další zlomky. Jeho úprava musí zachovat přípustnost všech vnitřních operací.
\end{definiceBox}


Buď nejprve uprav horní a dolní část a potom děl násobením převrácenou hodnotou, nebo rozšiř celý složený zlomek společným nenulovým násobkem vnitřních jmenovatelů. V obou postupech zachovej původní podmínky.

\blok{Řešené příklady}

\begin{prikladReseniBox}[title={Příklad 1},breakable=false]
Uprav $\dfrac{\dfrac{x}{y}}{\dfrac{2x}{3y}}$.
\tcblower
Vnitřní zlomky vyžadují $y\ne0$, celý dělitel navíc $x\ne0$. Potom $\frac{x}{y}\cdot\frac{3y}{2x}=\frac32$ při $x,y\ne0$.
\end{prikladReseniBox}


\begin{prikladReseniBox}[title={Příklad 2},breakable=false]
Uprav $\dfrac{\dfrac{x+1}{x-1}-1}{\dfrac{x+1}{x-1}+1}$.
\tcblower
Nejprve $x\ne1$ kvůli vnitřním jmenovatelům. Každou část upravíme zvlášť:
\begin{align*}
\frac{x+1}{x-1}-1&=\frac{(x+1)-(x-1)}{x-1}=\frac2{x-1},\\
\frac{x+1}{x-1}+1&=\frac{(x+1)+(x-1)}{x-1}=\frac{2x}{x-1}.
\end{align*}
Celá dolní část musí být nenulová, proto navíc $x\ne0$. Teprve pak dělíme:
\[
\frac{2/(x-1)}{2x/(x-1)}=\frac2{x-1}\cdot\frac{x-1}{2x}=\frac1x,
\qquad x\ne0,1.
\]
Krátíme součin, nikoli sčítance; oba původní zákazy ponecháme.
\end{prikladReseniBox}


\blok{Samostatné procvičení}

\begin{zadaniReseniBox}[title={Procvičení},breakable=false]
\begin{enumerate}[label=\alph*)]
\item Uprav $\dfrac{1/a}{2/a}$.
\item Uprav $\dfrac{1/x+1/(x+1)}{1/[x(x+1)]}$.
\item Uprav $\dfrac{(x^2-1)/(x-1)}{(x+1)/(x+2)}$.
\end{enumerate}
\tcblower
\begin{enumerate}[label=\alph*)]
\item $1/2$ při $a\ne0$; obě vnitřní části existují a dolní není nulová.
\item Pro $x\ne0,-1$ je dolní část nenulová. Rozšířením celého zlomku výrazem $x(x+1)$ dostaneme $2x+1$.
\item Původní zápis vyžaduje $x\ne1,-2$ a nenulový dolní zlomek navíc $x\ne-1$. Po krácení je výsledek $x+2$ při všech třech zákazech.
\end{enumerate}
\end{zadaniReseniBox}


\poznamka{Nech žáka nahlas označit hlavní zlomkovou čáru a celý dělitel. Podmínky dolní části nestačí pouze opsat z jejích jmenovatelů: musí být i nenulová. Rezervu se dvěma proměnnými lze dokončit doma.}


\blok{Navazující procvičování ve třídě}

\begin{zadaniReseniBox}[title={Rezerva: dvě cesty k témuž výsledku},breakable=false]
\begin{enumerate}[label=\alph*)]
\item Uprav $\dfrac{1+1/x}{1-1/x}$ dvěma způsoby: úpravou každé části a rozšířením celého zlomku.
\item Uprav $\dfrac{1/(x-1)-1/(x+1)}{1/(x-1)+1/(x+1)}$.
\end{enumerate}
\tcblower
\begin{enumerate}[label=\alph*)]
\item $x\ne0$ kvůli vnitřním zlomkům, $x\ne1$ kvůli dolní části. První cesta: $\frac{(x+1)/x}{(x-1)/x}=\frac{x+1}{x-1}$. Druhá: rozšířit celý zlomek nenulovým $x$, což dá tentýž tvar s oběma zákazy.
\item $x\ne\pm1$; horní část $2/(x^2-1)$, dolní $2x/(x^2-1)$, proto navíc $x\ne0$. Výsledek $1/x$ při $x\ne-1,0,1$.
\end{enumerate}
\end{zadaniReseniBox}

\blok{Volitelná rezerva}

\begin{prikladReseniBox}[title={Rozšíření procvičení},breakable=false]
Uprav $\dfrac{1/x+1/y}{1/x-1/y}$.
\tcblower
Nejprve $x,y\ne0$. Dolní část je $(y-x)/(xy)$, proto navíc $y\ne x$. Rozšířením výrazem $xy$ dostaneme $(x+y)/(y-x)$ na tomto oboru.
\end{prikladReseniBox}



\poznamka{Základ: původní P1 a P2 s doplněnými mezikroky, pak původní samostatné úlohy. Nová dvojice i původní dvě proměnné jsou rezervy, z nichž vybírej podle tempa. Domácí a) až c) jsou pro všechny, d) dobrovolně. U každého složeného zlomku žák musí ukázat, která čára je hlavní, a vypsat tři zdroje podmínek: vnitřní jmenovatele nahoře, dole a nenulovost celé dolní části.}

\blok{Domácí procvičení}

\begin{zadaniReseniBox}[title={Domácí úlohy},breakable=false]
\begin{enumerate}[label=\alph*)]
\item Uprav $\dfrac{2/x}{3/x}$.
\item Uprav $\dfrac{1/x+1}{1/x-1}$.
\item Uprav $\dfrac{1/(x-2)}{1/(x+2)}$.
\item Rezerva: uprav $\dfrac{x/y-y/x}{1/y}$.
\end{enumerate}
\tcblower
\begin{enumerate}[label=\alph*)]
\item $2/3$, $x\ne0$.
\item $\frac{x+1}{1-x}$, $x\ne0,1$.
\item $\frac{x+2}{x-2}$, $x\ne\pm2$.
\item $\frac{x^2-y^2}{x}$, $x,y\ne0$. Dolní část $1/y$ je na tomto oboru nenulová.
\end{enumerate}
\end{zadaniReseniBox}

\end{document}
